% Options for packages loaded elsewhere
\PassOptionsToPackage{unicode}{hyperref}
\PassOptionsToPackage{hyphens}{url}
%
\documentclass[
]{article}
\usepackage{amsmath,amssymb}
\usepackage{iftex}
\ifPDFTeX
  \usepackage[T1]{fontenc}
  \usepackage[utf8]{inputenc}
  \usepackage{textcomp} % provide euro and other symbols
\else % if luatex or xetex
  \usepackage{unicode-math} % this also loads fontspec
  \defaultfontfeatures{Scale=MatchLowercase}
  \defaultfontfeatures[\rmfamily]{Ligatures=TeX,Scale=1}
\fi
\usepackage{lmodern}
\ifPDFTeX\else
  % xetex/luatex font selection
\fi
% Use upquote if available, for straight quotes in verbatim environments
\IfFileExists{upquote.sty}{\usepackage{upquote}}{}
\IfFileExists{microtype.sty}{% use microtype if available
  \usepackage[]{microtype}
  \UseMicrotypeSet[protrusion]{basicmath} % disable protrusion for tt fonts
}{}
\makeatletter
\@ifundefined{KOMAClassName}{% if non-KOMA class
  \IfFileExists{parskip.sty}{%
    \usepackage{parskip}
  }{% else
    \setlength{\parindent}{0pt}
    \setlength{\parskip}{6pt plus 2pt minus 1pt}}
}{% if KOMA class
  \KOMAoptions{parskip=half}}
\makeatother
\usepackage{xcolor}
\usepackage{longtable,booktabs,array}
\usepackage{calc} % for calculating minipage widths
% Correct order of tables after \paragraph or \subparagraph
\usepackage{etoolbox}
\makeatletter
\patchcmd\longtable{\par}{\if@noskipsec\mbox{}\fi\par}{}{}
\makeatother
% Allow footnotes in longtable head/foot
\IfFileExists{footnotehyper.sty}{\usepackage{footnotehyper}}{\usepackage{footnote}}
\makesavenoteenv{longtable}
\setlength{\emergencystretch}{3em} % prevent overfull lines
\providecommand{\tightlist}{%
  \setlength{\itemsep}{0pt}\setlength{\parskip}{0pt}}
\setcounter{secnumdepth}{-\maxdimen} % remove section numbering
\ifLuaTeX
  \usepackage{selnolig}  % disable illegal ligatures
\fi
\IfFileExists{bookmark.sty}{\usepackage{bookmark}}{\usepackage{hyperref}}
\IfFileExists{xurl.sty}{\usepackage{xurl}}{} % add URL line breaks if available
\urlstyle{same}
\hypersetup{
  pdftitle={How far is cold fusion from known physics? A validated rate calculation, a two-sided bound on anomalous screening at rest, and a bound on the heat-only channel},
  pdfauthor={{[}Author names and affiliations to be added{]}},
  hidelinks,
  pdfcreator={LaTeX via pandoc}}

\title{How far is cold fusion from known physics? A validated rate
calculation, a two-sided bound on anomalous screening at rest, and a
bound on the heat-only channel}
\author{{[}Author names and affiliations to be added{]}}
\date{Draft, October 2026}

\begin{document}
\maketitle
\begin{abstract}
We compute deuterium--deuterium (and other hydrogen-isotope) fusion
rates for nuclei at rest with a log-grid Numerov solver for the
relative-motion ground state, \(\Gamma = A|\psi(0)|^2\). The solver
reproduces the D\(_2\), HD and H\(_2\) binding energies to \(\le 0.001\)
eV and the Koonin--Nauenberg molecular rates (d+d: \(2.5\times10^{-64}\)
vs \(3\times10^{-64}\) s\(^{-1}\); p+d: \(10^{-55.2}\) vs
\(\sim10^{-55}\) s\(^{-1}\)). For deuterons at the PdD lattice spacing
the rate is \(\sim10^{-124}\) s\(^{-1}\), about 114 orders of magnitude
below the \(\sim6\times10^{-11}\) s\(^{-1}\) per pair needed for 1 W
cm\(^{-3}\). Compressing a pair, every hydrogen-isotope fuel except p+p
reaches that level at a common spacing of 7--8 pm. We then model the
anomalous electron screening reported in low-energy beam experiments on
deuterated metals (165--800 eV) as a profile of strength fixed by the
beam data and unknown range \(\ell\). As \(\ell\) grows, the rate jumps
once the extra attraction binds a new short-range d--d state. Combining
the beam constraint with the null neutron limit for loaded Pd and Ti at
rest (\(<2\times10^{-26}\) s\(^{-1}\) per pair) bounds \(\ell\) to
\(\approx 0.01\)--\(0.3\) Å across 40+ model variants. In every variant,
rates approaching 1 W cm\(^{-3}\) occur only at ranges already excluded
by the neutron limit. Finally, we bound the ``heat-only'' channel, in
which fusion energy goes directly to the lattice. Suppressing neutrons
by even \(10^{-3}\) requires a lattice decay width exceeding the 23.85
MeV transition energy. Electromagnetic, multi-phonon and coherent
many-nucleus mechanisms fall short by \$\sim\$11, \(\sim10^9\) and
\(\gtrsim 12\) orders of magnitude, and coherent transfer is
additionally excluded by causality. We propose three experiments that
test the remaining assumptions: a 5.5 MeV gamma search for p+d fusion in
H/D co-loaded Pd at rest, beam screening in a heavy-fermion compound
with a temperature switch, and same-target p+d vs d+d screening.
\end{abstract}

\hypertarget{introduction}{%
\section{1. Introduction}\label{introduction}}

Claims of nuclear fusion in deuterium-loaded metals at room temperature
{[}Fleischmann--Pons 1989; Jones et al.~1989{]} were followed within
months by calculations of the molecular d+d fusion rate {[}Koonin \&
Nauenberg 1989{]}, by general bounds on barrier penetration in many-body
systems at equilibrium {[}Leggett \& Baym 1989{]}, and by null neutron
searches {[}e.g.~Aberdam et al.~1990{]}. A large Google-funded
re-examination found no evidence for the effect {[}Berlinguette et
al.~2019{]}.

One genuine anomaly has since been established by accelerator
experiments. Low-energy d(d,p)t cross-sections in deuterated metals
imply \textbf{electron screening energies} \(U_e \approx 200\)--\(800\)
eV {[}Raiola et al.~2002, 2004; Czerski et al.{]}, an order of magnitude
above the \(\approx 25\) eV measured with gas targets and the
\(\approx 28\) eV adiabatic limit. Recent work reports a 15\% increase
in beam-driven fusion with electrochemical loading {[}Chen et
al.~2025{]} and a sub-keV yield plateau in Pd and Ti deuterides
{[}Karahadian et al.~2026{]}. Whether such effects persist for nuclei
\emph{at rest}, and how far they could go towards practical rates, has
not been quantified jointly with the null searches.

This paper asks a narrow question: \textbf{how far, in orders of
magnitude, is room-temperature fusion in metals from known physics, and
which variables can close the gap?} We use one validated solver
throughout. All code, notebooks and figures are public {[}repository{]}.

\hypertarget{method}{%
\section{2. Method}\label{method}}

For a bound pair, the fusion rate is
\[\Gamma = A\,|\psi(0)|^2, \qquad A = \frac{S(0)\,c}{\pi\alpha\,\mu c^2},\]
where \(S(0)\) is the zero-energy astrophysical S-factor and \(\psi\)
the s-wave ground state of the relative motion in potential \(V(R)\).
For d+d with \(S(0) = 109.3\) keV b (Bosch--Hale, both branches),
\(A = 1.52\times10^{-16}\) cm\(^3\) s\(^{-1}\). We solve
\(u'' = 2\mu(V-E)u\) on a logarithmic grid from 0.5 fm, starting from
the regular repulsive-Coulomb solution \(u\simeq R(1+\mu R)\). The
energy comes from node-counting bisection, and amplitudes are tracked in
log space (\(\psi(0)\) is \(\sim10^{-37}\) of its peak).

\textbf{Potentials.} \emph{Molecular:} the H\(_2\) Born--Oppenheimer
curve {[}Kolos \& Wolniewicz 1965{]}, splined to the He united-atom
limit; heavier screening particles (mass \(m^*\)) via the exact scaling
\(V_{m^*}(R)=m^*V_1(m^*R)\). \emph{Lattice:} Thomas--Fermi-screened
repulsion (0.57 Å for Pd) plus a harmonic restoring force on the
separation set by the PdD optical phonon (37 meV). \emph{Anomalous
screening:} \(\Delta V(R) = -\Delta U\,g(R/\ell)\) with \(g\)
exponential, Gaussian, step or Yukawa-type. \(\Delta U\) is fixed by
requiring \(\Delta V = -U_{\rm meas}\) at the beam turning point
(\(\approx 0.004\) Å). When \(\Delta V\) creates a bound state at short
range, we take the effective rate as the smaller of the bound-state
fusion rate and the WKB tunneling rate into it.

\textbf{Validation} (Table 1). The solver reproduces D\(_2\), HD and
H\(_2\) binding energies to \(\le 0.001\) eV, the D\(_2\) zero-point
energy to \(<1\%\), and the Koonin--Nauenberg rates for d+d and p+d.~It
is converged to 0.002 decades in step size and starting radius. The main
systematic is the short-range potential: a 0.05 Ha error below 0.3 Å
shifts the rate by 0.7 decades. We adopt \(\pm1\) decade as the working
uncertainty.

\begin{longtable}[]{@{}
  >{\raggedright\arraybackslash}p{(\columnwidth - 4\tabcolsep) * \real{0.3333}}
  >{\raggedright\arraybackslash}p{(\columnwidth - 4\tabcolsep) * \real{0.3333}}
  >{\raggedright\arraybackslash}p{(\columnwidth - 4\tabcolsep) * \real{0.3333}}@{}}
\toprule\noalign{}
\begin{minipage}[b]{\linewidth}\raggedright
Quantity
\end{minipage} & \begin{minipage}[b]{\linewidth}\raggedright
This work
\end{minipage} & \begin{minipage}[b]{\linewidth}\raggedright
Reference
\end{minipage} \\
\midrule\noalign{}
\endhead
\bottomrule\noalign{}
\endlastfoot
D\(_2\) binding \(D_0\) & 4.5556 eV & 4.5562 eV \\
HD / H\(_2\) binding & 4.5129 / 4.4771 eV & 4.5139 / 4.4781 eV \\
d+d rate in D\(_2\) & \(2.5\times10^{-64}\) s\(^{-1}\) &
\(3\times10^{-64}\) {[}KN89{]}; \(5.7\times10^{-64}\) {[}2025{]} \\
p+d rate in HD & \(10^{-55.2}\) s\(^{-1}\) & \(\sim10^{-55}\)
{[}KN89{]} \\
Thermal reactivity (Gamow integral) & within 1--2\% & Bosch--Hale fit,
0.2--50 keV \\
\end{longtable}

\hypertarget{results}{%
\section{3. Results}\label{results}}

\hypertarget{the-gap}{%
\subsection{3.1 The gap}\label{the-gap}}

1 W cm\(^{-3}\) in PdD requires \(\sim6\times10^{-11}\) fusions per pair
per second. A d--d pair at the lattice spacing (2.75 Å) gives
\(10^{-124}\) s\(^{-1}\), a gap of \textbf{\textasciitilde114 orders of
magnitude}. The often-quoted \textasciitilde53 orders is against a
D\(_2\) molecule (0.74 Å), which the lattice does not provide. Varying
spacing (1.5--2.75 Å), phonon energy and Thomas--Fermi length moves the
lattice rate by at most tens of orders.

\hypertarget{closer-and-lighter}{%
\subsection{3.2 Closer and lighter}\label{closer-and-lighter}}

Scaling the molecular potential by a heavier screening particle
compresses the pair. Watt-level rates need \(m^*\approx 9.2\) for d+d
(Koonin \& Nauenberg: 5--10), i.e.~a \textbf{pair spacing of 8.0 pm}.
Repeating this for all hydrogen-isotope pairs, p+d, p+t, d+d and d+t all
reach watt level at \textbf{7.3--8.2 pm}; p+p never does
(weak-interaction S-factor). Lighter pairs lead at long range (p+d is
\(10^{8}\) above d+d in a molecule, \(10^{17}\) at lattice spacing), but
larger S-factors win near the watt zone. An independent calculation for
a lattice-type pair held at fixed spacing by an external force gives the
same 7.3 pm, at an energy cost of \textasciitilde170 eV per pair
(Boltzmann factor \(10^{-2900}\) at 300 K). Deuterium clusters in a Pd
vacancy (N = 2--6) do not help: mutual repulsion places the nearest pair
at 2.8--3.8 Å.

\hypertarget{anomalous-screening-at-rest-a-two-sided-bound}{%
\subsection{3.3 Anomalous screening at rest: a two-sided
bound}\label{anomalous-screening-at-rest-a-two-sided-bound}}

Fig. 1 shows the effective rate for a lattice pair along the
beam-consistent curve. If the screening reaches only the few-hundredths
of an ångström the beam data probe, the rate stays below \(10^{-100}\)
s\(^{-1}\). As \(\ell\) increases, the rate \textbf{jumps} when
\(\Delta V\) binds a new d--d state at 0.07--0.24 Å (deeper and closer
for larger \(U_{\rm meas}\)). It is then limited by tunneling into that
state and by fusion within it.

The null neutron limit for Pd and Ti at rest, \(<2\times10^{-26}\)
s\(^{-1}\) per pair {[}Aberdam et al.~1990{]}, then bounds \(\ell\). For
an exponential profile, \(\ell < 0.154\) Å (300 eV), \(0.137\) Å (575
eV) and \(0.129\) Å (800 eV). Across 40 variants (profile shape,
harmonic vs barrier-capped lattice, spacing, Thomas--Fermi length,
phonon energy, attempt frequency \(\times10^{\pm3}\), weaker null limit)
plus the 800 eV case, the bound spans \textbf{0.012--0.315 Å} (median
0.137 Å). It is dominated by profile shape and screening length; lattice
stiffness and tunneling kinetics hardly matter. At 165 eV and for step
profiles the static limit does not constrain \(\ell\), but those cases
never come within 20 orders of watt level.

\textbf{Watts imply neutrons.} Several variants reach \(\ge 1\) W
cm\(^{-3}\) at \(U_{\rm meas}\ge 575\) eV. In every one, the rate
crosses the neutron limit at a \emph{shorter} range than it reaches watt
level (e.g.~Yukawa profile at 800 eV: 0.012 Å vs 0.016 Å). No variant
jumps from below the limit to watt level.

\hypertarget{pd-in-mixed-hd-lattices}{%
\subsection{3.4 p+d in mixed H/D
lattices}\label{pd-in-mixed-hd-lattices}}

p+d → \(^3\)He + γ (5.49 MeV) has no neutron branch. With the same
beam-consistent screening, the p+d lattice rate peaks at \(10^{-12.3}\)
s\(^{-1}\) (575 eV), so even at a 50/50 H/D mix it is \textasciitilde2.5
orders short of watt level. Past the bound-state jump, d+d and p+d
converge to within an order of magnitude, so the coexisting d+d pairs
would still violate the neutron limit first. p+d also cannot be
``heat-only'': 1 W would emit \(\sim10^{12}\) gammas per second.

\hypertarget{the-heat-only-channel}{%
\subsection{3.5 The heat-only channel}\label{the-heat-only-channel}}

The d+d compound state \(^4\)He\(^*\) (23.85 MeV) decays strongly with
widths \(\Gamma_{\rm strong}\sim0.5\)--6 MeV (lifetimes
\(10^{-22}\)--\(10^{-21}\) s), and radiatively with
\(\Gamma_\gamma/\Gamma_p\approx 10^{-7}\) {[}Zhou et al.~2006{]}.
Suppressing neutrons by \(s_n\) needs a lattice width
\(\Gamma_{\rm lat}\gtrsim\Gamma_{\rm strong}/s_n\), which is
\textbf{already 20--260× the transition energy at \(s_n=10^{-3}\)} and
\(10^{10}\)--\(10^{11}\)× at \(10^{-12}\). A state that wide is not a
state, so any heat-only process cannot go through \(^4\)He\(^*\). For
mechanisms that do:

\begin{itemize}
\tightlist
\item
  \textbf{electromagnetic} (internal conversion, pair creation) falls
  short by \textasciitilde11--20 orders;
\item
  \textbf{direct multi-phonon} emission (Englman--Jortner energy-gap
  law, \(6.4\times10^8\) phonons) gives rates \(\sim10^{-4\times10^9}\)
  s\(^{-1}\) for any plausible coupling;
\item
  \textbf{coherent many-nucleus} coupling needs
  \(\sim10^{15}\)--\(10^{33}\) phase-locked nuclei (vs \(\sim10^3\) in
  coherent domains).
\end{itemize}

It is also \textbf{causally excluded}: in the compound-nucleus lifetime,
light travels \(<0.003\) Å, far less than the distance to the nearest
neighbouring deuteron.

\hypertarget{other-levers}{%
\subsection{3.6 Other levers}\label{other-levers}}

Cold metallic deuterium reaches 1 W cm\(^{-3}\) pycnonuclearly only at
\(\sim1.6\times10^4\) g cm\(^{-3}\). Coherent resonant tunneling through
an intermediate well adds nothing beyond sequential tunneling once
energies are spread over more than the resonance width (≲\(10^{-31}\) eV
vs \(\sim10^{-2}\) eV). Muon-catalysed d--t fusion, the one
room-temperature route that works, is \textasciitilde1.5× from
scientific breakeven and \textasciitilde9× from net electricity, a gap
that depends on sticking and cycling measurements now being made at high
density (MuFusE).

\hypertarget{proposed-experiments}{%
\section{4. Proposed experiments}\label{proposed-experiments}}

\textbf{E1: p+d gamma search at rest.} A 10 g Pd foil co-loaded with H
and D (0.45/0.45), steady and cycled loading, a shielded HPGe/LaBr\(_3\)
detector with muon veto, run for 30 days, with H-only, D-only and
unloaded controls. With assumed backgrounds this reaches
\(\sim10^{-26}\) (surface) and \(4\times10^{-28}\) (underground) per
p--d pair per s, extending the range bound to a neutron-free channel.

\textbf{E2: heavy-fermion screening with a temperature switch.} d(d,p)t
at \(E_d\) = 5--30 keV on a deuterium-implanted heavy-fermion compound
(Ce- or Yb-based) and its non-f analogue, at 4 K and 300 K, with in-situ
ERDA depth profiling. ``Heavy-electron'' screening predicts a large rise
only below the coherence temperature; standard theory predicts no
change. About 6 h per point at 10 μA.

\textbf{E3: same-target p+d vs d+d screening.} Alternating
H\(^+\)/D\(^+\) beams matched in stopping range on one deuterated
target. This separates true screening (equal \(U_e\)), depth-profile
artifacts (range-dependent) and a d+d-specific nuclear resonance
(\(U_e^{dd}>U_e^{pd}\)). A 5σ answer needs \textasciitilde2.6\% p+d
statistics, about 1.5 days per point at 1 mA.

\hypertarget{discussion-and-limitations}{%
\section{5. Discussion and
limitations}\label{discussion-and-limitations}}

\textbf{Relation to prior work.} Leggett \& Baym {[}1989{]} bounded
barrier penetration for any many-body system at equilibrium; our
static-screening conclusion agrees. The contribution here is to use the
\emph{measured} beam screening, model its unknown range, and confront it
with the \emph{null} neutron data, which gives a numerical bound on the
range at rest and the ``watts imply neutrons'' result across model
variants. Czerski, Heide \& Huke {[}2006{]} extrapolated effective
screening energies to cold-fusion energies and compared with positive
claims. Their full text should be compared in detail before any novelty
claim is made.

\textbf{Limitations.} (i) The short-range H\(_2\) potential uses
recalled Kolos--Wolniewicz values. They are supported by the
binding-energy and Koonin--Nauenberg checks but were not re-read from
the original table. (ii) Several literature inputs were verified only at
abstract level (see VERIFICATION.md). (iii) The lattice model is a
single radial pair; anisotropy, many-body exchange and defects beyond a
spherical vacancy are not modelled. (iv) The screening profiles are
phenomenological. (v) The heat-only analysis uses order-of-magnitude
internal-conversion and pair-creation coefficients and the weak-coupling
energy-gap law.

\hypertarget{conclusions}{%
\section{6. Conclusions}\label{conclusions}}

Within standard physics and the stated assumptions:

\begin{enumerate}
\def\labelenumi{\arabic{enumi}.}
\tightlist
\item
  Room-temperature fusion in metals is \textasciitilde114 orders of
  magnitude from 1 W cm\(^{-3}\).
\item
  Every fuel needs pairs held at \textasciitilde7--8 pm, costing
  \textasciitilde170 eV per pair.
\item
  The anomalous screening seen in beam experiments cannot extend beyond
  \textasciitilde0.01--0.3 Å at rest without having produced neutrons
  already excluded in 1990, and any configuration approaching watts
  would have produced them.
\item
  A heat-only channel cannot dispose of the energy of any process that
  forms excited \(^4\)He.
\end{enumerate}

The remaining open questions are experimental and well posed (E1--E3).

\textbf{Data and code availability.} All notebooks, the solver
(\texttt{cfsolve.py}), the claim-checking tool and figures:
https://github.com/m1110/cold-fusion-gap

\textbf{Acknowledgements.} Analysis code and drafting were assisted by
an AI model (Claude, Anthropic); all results are reproducible from the
repository.

\hypertarget{references-to-be-completed-and-checked-against-primary-sources}{%
\section{References (to be completed and checked against primary
sources)}\label{references-to-be-completed-and-checked-against-primary-sources}}

\begin{itemize}
\tightlist
\item
  Aberdam D. et al., Phys. Rev.~Lett. 65, 1196 (1990).
\item
  Berlinguette C. P. et al., Nature 570, 45 (2019).
\item
  Bosch H.-S. \& Hale G. M., Nucl. Fusion 32, 611 (1992) (and erratum).
\item
  Chen K.-Y. et al., Nature 644, 640 (2025).
\item
  Ciuchi S. et al., Eur. Phys. J. C 72, 2193 (2012).
\item
  Czerski K., Heide P. \& Huke A., in Condensed Matter Nuclear Science
  (ICCF11), 228 (2006).
\item
  Czerski K. et al., EPL 113, 22001 (2016).
\item
  Englman R. \& Jortner J., Mol. Phys. 18, 145 (1970).
\item
  Huke A., Czerski K. \& Heide P., arXiv:0805.4538 (2008).
\item
  Jones S. E. et al., Nature 338, 737 (1989).
\item
  Karahadian M. et al., Nat. Commun. 17, 8845 (2026).
\item
  Kolos W. \& Wolniewicz L., J. Chem. Phys. 43, 2429 (1965).
\item
  Koonin S. E. \& Nauenberg M., Nature 339, 690 (1989).
\item
  Leggett A. J. \& Baym G., Phys. Rev.~Lett. 63, 191 (1989); Nature 340,
  45 (1989).
\item
  Raiola F. et al., Eur. Phys. J. A 13, 377 (2002); 19, 283 (2004).
\item
  Tilley D. R., Weller H. R. \& Hale G. M., Nucl. Phys. A 541, 1 (1992).
\item
  Zhou J. et al., Chin. Phys. Lett. (2006), d(d,γ)⁴He branching at 20
  keV.
\end{itemize}

\end{document}
